%----------------------------------------------------------------------------------------
% INTEGRATIEMETHODEN. Alle opgaven zijn eigen werk.
%
% Opbouw: Start, leerstof, en op een nieuwe pagina de Opdrachten, Koppel en Test jezelf.
% Labels van dit hoofdstuk beginnen met h3-.
%----------------------------------------------------------------------------------------

\mychapter{3}{Integratiemethoden}

\lkr{Leerplan B+S': LPD 25, met onmiddellijke integratie, splitsing, eenvoudige substitutie en partiële integratie. Splitsen in partieelbreuken en andere substituties staan bij de suggesties voor het complementair gedeelte, en dus niet in dit hoofdstuk. De grootste moeilijkheid is niet de techniek zelf maar de keuze: daarom sluit de leerstof af met een stuk over welke methode je kiest, en staan de opdrachten door elkaar. Laat leerlingen hun antwoord telkens controleren met de afgeleide. Vermeld ook dat veel integralen geen primitieve in gekende functies hebben (leerplan), en dat ICT dan een waarde kan geven.}

Dankzij de hoofdstelling is een integraal berekenen hetzelfde als een primitieve zoeken. Maar afleiden en integreren zijn niet even gemakkelijk. Voor afleiden heb je regels voor elk product en elke samengestelde functie. Omgekeerd bestaan die niet: de integraal van een product is niet het product van de integralen.

Toch geeft elke regel om af te leiden, achterstevoren gelezen, een methode om te integreren. De kettingregel wordt substitutie, de productregel wordt partiële integratie. Die twee methodes en het splitsen van een integrand vormen de kern van dit hoofdstuk. De moeilijkheid zit minder in het rekenen dan in het kiezen: daarom sluiten we af met de vraag welke methode je wanneer gebruikt.

\begin{oefeningen}[Start]
\oef{h3-start-1}

\oef{h3-verken}
\end{oefeningen}

\lkr{Oefening \ref{oef:h3-verken} toont de twee methodes van dit hoofdstuk als omgekeerde regels: de kettingregel en de productregel. Laat vraag f gerust even mislukken.}

%========================================================================================
\section{De onbepaalde integraal}

Voor een bepaalde integraal had je één primitieve nodig. Soms wil je ze allemaal: als een extra voorwaarde moet zeggen welke de juiste is, of om een methode rustig op te schrijven zonder telkens grenzen mee te dragen. Daarvoor bestaat een aparte notatie, zonder grenzen.

\begin{theorie}
\begin{stelling}{Onbepaalde integraal}
De \textbf{\textcolor{accent}{onbepaalde integraal}} van $f$ is de verzameling van alle primitieven van $f$ op een interval:

\vspace{1mm}
\qquad $\formule{\int f(x)\dx = F(x) + C}$ \quad met $F'(x) = f(x)$
\end{stelling}

De \textbf{\textcolor{accent}{lineariteit}} uit hoofdstuk 1 blijft gelden: \quad $\formule{\int \bigl(r f(x) + s g(x)\bigr)\dx = r\int f(x)\dx + s\int g(x)\dx}$
\end{theorie}

De bepaalde integraal is een getal, de onbepaalde een verzameling functies. Vergeet de $C$ dus niet, en controleer met de afgeleide. Alle basisprimitieven staan in de tabel van hoofdstuk 2. Rekenen met die tabel heet \emph{onmiddellijke integratie}.

\voorbeeld{}{Bereken $\ig (6x^2 - 4x + 3)\dx$.\label{vb:h3-onb}}{%
\uitleg{Er is geen voorwaarde, dus de uitkomst is een hele verzameling functies, met een willekeurige constante.}
\begin{stappen}
\ig (6x^2 - 4x + 3)\dx = 2x^3 - 2x^2 + 3x + C & Per term een primitieve. Elke term zou een constante geven, maar samen is dat één $C$.
\end{stappen}}

Soms wil je niet de hele verzameling, maar precies één primitieve. Dan is er een extra voorwaarde nodig.

\voorbeeld{}{Bepaal de primitieve $F$ van $f(x) = 3x^2 - \dfrac{1}{x^2}$ met $F(1) = 5$.\label{vb:h3-voorwaarde}}{%
\uitleg{Eerst zoeken we alle primitieven, met de constante erbij. De voorwaarde bepaalt daarna welke constante we nodig hebben.}
\naast[]{\begin{stappenvrij}{7.2cm}{3.9cm}
F(x) = x^3 + \dfrac{1}{x} + C & $-\dfrac{1}{x^2}$ is $-x^{-2}$, met primitieve $x^{-1}$. \\[2ex]
1 + 1 + C = 5 \ \Rightarrow \ C = 3 & Vul de voorwaarde in. Dus $F(x) = x^3 + \tfrac1x + 3$.
\end{stappenvrij}}{\begin{tikzpicture}[x=1cm,y=.26cm,font=\scriptsize]
  \rooster{0}{1}{2}{0}{2}{8}
  \foreach \c in {1,5}{\kromme{.5}{1.3}{(\x)^3+1/\x+\c}}
  \kromme[very thick]{.5}{1.3}{(\x)^3+1/\x+3}
  \fill[accent] (1,5) circle (2pt);
  \node[right] at (1,4.6) {$(1, 5)$};
\end{tikzpicture}}}

De figuur toont wat de voorwaarde doet: uit de hele familie krommen kiest ze de ene die door het punt $(1, 5)$ gaat.

\nujij{Bereken $\int (x^2 - \cos x)\dx$. Bepaal de primitieve van $f(x) = 2x + 1$ met $F(0) = 7$.}{$\tfrac{x^3}{3} - \sin x + C$, en $F(x) = x^2 + x + 7$.}

\begin{oefeningen}
\oef{h3-onb}
\end{oefeningen}

%========================================================================================
\section{Integratie door splitsing}

Niet elke integrand staat meteen in de tabel, maar vaak is het een som in vermomming: een breuk met één term in de noemer, een haakje dat je kunt uitwerken of een goniometrische uitdrukking. Eens uitgeschreven valt ze uiteen in termen die je kent. Dat heb je in hoofdstuk 2 al gedaan bij bepaalde integralen.

\voorbeeld{}{Bereken $\ig \dfrac{x^2 - 3x + 2}{x}\dx$.\label{vb:h3-split}}{%
\uitleg{De integrand staat niet in de tabel, maar de noemer is één term. Dan kunnen we de breuk splitsen.}
\begin{stappen}
\ig \Bigl(x - 3 + \dfrac{2}{x}\Bigr)\dx = \dfrac{x^2}{2} - 3x + 2\ln |x| + C & We delen elke term van de teller door $x$. Let op: $\tfrac2x$ geeft een logaritme.
\end{stappen}}

\voorbeeld{}{Bereken $\ig \tan^2 x\dx$.\label{vb:h3-tan}}{%
\uitleg{Ook $\tan^2 x$ staat niet in de tabel, maar de afgeleide van $\tan x$ wel. Een goniometrische formule brengt die in beeld.}
\begin{stappen}
\tan^2 x = \dfrac{1}{\cos^2 x} - 1 & Deel $\sin^2 x + \cos^2 x = 1$ door $\cos^2 x$. \\[2.5ex]
\ig \Bigl(\dfrac{1}{\cos^2 x} - 1\Bigr)\dx = \tan x - x + C & Zo krijg je twee termen uit de tabel.
\end{stappen}}

\begin{werkwijze}
Schrijf de integrand als een som van termen die in de tabel staan. Soms gaat dat door uit te werken of een breuk te splitsen, soms met een goniometrische formule zoals $\tan^2 x = \dfrac{1}{\cos^2 x} - 1$.
\end{werkwijze}

\nujij{Bereken $\int \dfrac{(x - 1)^2}{x}\dx$.}{$\int (x - 2 + \tfrac1x)\dx = \tfrac{x^2}{2} - 2x + \ln |x| + C$.}

\begin{oefeningen}
\oef{h3-split}
\end{oefeningen}

%========================================================================================
\section{Integratie door substitutie}

\subsection{De substitutieregel}

Hoe integreer je $x\cos(x^2)$? Voor een product is er geen regel. Maar leid je $\sin(x^2)$ af, dan krijg je $\cos(x^2) \cdot 2x$: zo'n product ontstaat precies bij de kettingregel, met de afgeleide van de binnenste functie als extra factor. De substitutieregel leest de kettingregel achterstevoren. Je geeft de binnenste functie een nieuwe naam, $u$, en de integraal wordt een stuk eenvoudiger.


\begin{theorie}
\begin{stelling}{Substitutieregel}
Is $F$ een primitieve van $f$, en stel $u = g(x)$, dan is $\mathrm{d}u = g'(x)\dx$ en

\vspace{1mm}
\qquad $\formule{\int f(g(x))\,g'(x)\dx = \int f(u)\du = F(g(x)) + C}$

\vspace{1.5mm}
Bij een bepaalde integraal gaan ook de grenzen mee:

\vspace{1mm}
\qquad $\formule{\int_a^b f(g(x))\,g'(x)\dx = \int_{g(a)}^{g(b)} f(u)\du}$
\end{stelling}

\begin{bewijs}
Uit de kettingregel en $F' = f$ volgt
\bewijsregel{\bigl(F(g(x))\bigr)' = F'(g(x)) \cdot g'(x) = f(g(x)) \cdot g'(x) .}
Dus $F(g(x))$ is een primitieve van de integrand, en met $u = g(x)$ is
\bewijsregel{\ig f(g(x))\,g'(x)\dx = F(g(x)) + C = F(u) + C = \ig f(u)\du .}
Met grenzen geeft de hoofdstelling, twee keer toegepast,
\bewijsregel{\ig_a^b f(g(x))\,g'(x)\dx = F(g(b)) - F(g(a)) = \ig_{g(a)}^{g(b)} f(u)\du .}
\end{bewijs}
\end{theorie}

Waar komt die $\mathrm{d}u = g'(x)\dx$ vandaan? Uit $u = g(x)$ volgt $g'(x) = \dfrac{\mathrm{d}u}{\mathrm{d}x}$ in de notatie van Leibniz. Reken je daar informeel mee alsof het een breuk is, dan breng je $\mathrm{d}x$ naar de andere kant en krijg je $\mathrm{d}u = g'(x)\dx$. Dat is de kettingregel in differentiaalvorm. Een echte breuk is het niet, maar het bewijs hierboven toont dat je zo mag rekenen.

\voorbeeld{}{Bereken $\ig x \cos(x^2)\dx$.\label{vb:h3-sub}}{%
\uitleg{Bij elke substitutie stellen we dezelfde vraag: welke functie zit binnenin, en staat haar afgeleide ernaast? Hier staat $x^2$ in de cosinus, en de $x$ ernaast verraadt de afgeleide.}
\begin{stappen}
\text{Stel } u = x^2, \quad \mathrm{d}u = 2x\dx & De afgeleide van $x^2$ staat er, op de factor 2 na. Dus $x\dx = \tfrac12\du$. \\[1ex]
\ig x \cos(x^2)\dx = \tfrac12 \ig \cos u\du & Alles staat nu in $u$, en dat is een basisintegraal. \\[2.5ex]
\phantom{\ig x \cos(x^2)\dx} = \tfrac12 \sin(x^2) + C & Terug naar $x$. Controle: afleiden geeft de integrand.
\end{stappen}}

\voorbeeld{}{Bereken $\ig (3x - 1)^5\dx$.\label{vb:h3-sublin}}{%
\uitleg{De binnenste functie is hier lineair, dus haar afgeleide is een constante. Die breng je gewoon voor de integraal.}
\begin{stappenbreed}
\text{Stel } u = 3x - 1, \quad \mathrm{d}u = 3\dx & Dus $\dx = \tfrac13\du$. \\[1ex]
\ig (3x-1)^5\dx = \tfrac13 \ig u^5\du = \dfrac{(3x-1)^6}{18} + C & Uitwerken tot een veelterm van graad 5 kan ook, maar is veel meer werk.
\end{stappenbreed}}

\begin{werkwijze}
Zoek een samengestelde functie waarvan de binnenste functie $g(x)$ een afgeleide heeft die, op een constante na, als factor in de integrand staat. Stel $u = g(x)$ en bereken $\mathrm{d}u$. Schrijf alles met $u$: er mag geen $x$ meer overblijven. Integreer, en vervang $u$ weer door $g(x)$.
\end{werkwijze}

\nujij{Bereken $\int x^2 \e^{x^3}\dx$.}{$u = x^3$, $\mathrm{d}u = 3x^2\dx$: $\tfrac13 \e^{x^3} + C$.}

\begin{oefeningen}
\oef{h3-sub}
\end{oefeningen}

Substitutie werkt niet altijd: de afgeleide van de binnenste functie moet ook echt als factor in de integrand staan. Neem $\int \cos(x^2)\dx$, zonder de $x$ van daarnet. Iemand stelt $u = x^2$ en $\mathrm{d}u = 2x\dx$, dus $\dx = \dfrac{\mathrm{d}u}{2x}$. Dan staat er
\[ \int \cos(x^2)\dx = \int \frac{\cos u}{2x}\du , \]
met nog een $x$ in de integrand. Die mag er niet staan: je integreert naar $u$, en je kunt $\tfrac{1}{2x}$ niet zomaar voor de integraal brengen, want $x$ hangt af van $u$. Vervang je $x$ door $\sqrt{u}$, dan is de integraal ook niet eenvoudiger geworden. Niet alles is naar $u$ om te zetten, dus deze substitutie helpt niet. Het antwoord $\tfrac12\sin(x^2) + C$ is hier dan ook fout: afleiden geeft $x\cos(x^2)$.

\subsection{Vormen die vaak terugkomen}

Sommige substituties komen zo vaak voor dat je het resultaat beter meteen herkent. De eerste vorm kun je zelfs zien. De grafiek van $\sin 2x$ is die van $\sin x$, horizontaal samengedrukt met factor 2. De hoogtes blijven gelijk en de breedte halveert, dus ook de oppervlakte halveert: van 2 naar 1. Vandaar de factor $\tfrac{1}{a}$.

\figuren{\begin{tikzpicture}[x=.8cm,y=.9cm,font=\scriptsize]
  \gebied{0}{3.1416}{sin(\x r)}
  \roosterleeg{0}{1}{4}{0}{1}{1}
  \kromme{0}{3.1416}{sin(\x r)}
  \node[below] at (3.1416,0) {$\pi$};
  \node[left] at (0,1) {1};
  \node[text=accent] at (1.57,.4) {$2$};
  \node[above,text=accent] at (2.7,.7) {$\sin x$};
\end{tikzpicture}\qquad\qquad
\begin{tikzpicture}[x=.8cm,y=.9cm,font=\scriptsize]
  \gebied{0}{1.5708}{sin(2*\x r)}
  \roosterleeg{0}{1}{4}{0}{1}{1}
  \kromme{0}{1.5708}{sin(2*\x r)}
  \node[below] at (1.5708,0) {$\frac{\pi}{2}$};
  \node[left] at (0,1) {1};
  \node[text=accent] at (.785,.4) {$1$};
  \node[above,text=accent] at (2.1,.7) {$\sin 2x$};
\end{tikzpicture}}

Elk van de volgende formules kun je zelf nakijken met de afgeleide.

\begin{theorie}
\begin{stelling}{Standaardvormen}
{\centering\arrayrulecolor{accent}\renewcommand{\arraystretch}{2.4}%
\begin{tabular}{|>{\centering\arraybackslash}m{4.2cm}|>{\centering\arraybackslash$}m{7.6cm}<{$}|}
\hline
Lineaire binnenste functie & \formule{\int f(ax + b)\dx = \dfrac{1}{a}\,F(ax + b) + C} \\ \hline
Macht ($r \neq -1$) & \formule{\int g'(x)\,g(x)^r\dx = \dfrac{g(x)^{r+1}}{r+1} + C} \\ \hline
Afgeleide in de teller & \formule{\int \dfrac{g'(x)}{g(x)}\dx = \ln |g(x)| + C} \\ \hline
Exponentiële functie & \formule{\int g'(x)\,\e^{g(x)}\dx = \e^{g(x)} + C} \\ \hline
Met een constante $a > 0$ & \formule{\int \dfrac{1}{a^2 + x^2}\dx = \dfrac{1}{a}\bgtan\dfrac{x}{a} + C} \\ \hline
Met een constante $a > 0$ & \formule{\int \dfrac{1}{\sqrt{a^2 - x^2}}\dx = \bgsin\dfrac{x}{a} + C} \\ \hline
\end{tabular}\par}
\end{stelling}
\end{theorie}

\voorbeeld{}{Bereken $\ig \tan x\dx$, \ $\ig \dfrac{x}{x^2 + 1}\dx$ en $\ig \dfrac{1}{4 + x^2}\dx$.\label{vb:h3-vormen}}{%
\uitleg{Dit zijn drie vormen uit de tabel hierboven. We schrijven geen $u$ meer op, maar herkennen de vorm meteen. Let goed op het verschil tussen de laatste twee: er staat wel of geen $x$ in de teller.}
\begin{stappenbreed}
\ig \dfrac{\sin x}{\cos x}\dx = -\ig \dfrac{-\sin x}{\cos x}\dx = -\ln |\cos x| + C & Op een minteken na staat de afgeleide van $\cos x$ in de teller. \\[3ex]
\tfrac12 \ig \dfrac{2x}{x^2 + 1}\dx = \tfrac12 \ln(x^2 + 1) + C & De afgeleide in de teller; haakjes volstaan, want $x^2 + 1 > 0$. \\[2.5ex]
\ig \dfrac{1}{2^2 + x^2}\dx = \tfrac12 \bgtan \dfrac{x}{2} + C & Geen $x$ in de teller, dus geen logaritme. Hier is $a = 2$.
\end{stappenbreed}}

\nujij{Bereken $\int \dfrac{\cos x}{1 + \sin x}\dx$ en $\int \dfrac{1}{\sqrt{9 - x^2}}\dx$.}{$\ln(1 + \sin x) + C$ en $\bgsin\tfrac{x}{3} + C$.}

\voorbeeld{}{Voor elke vorm van de tabel nog een integraal. (Niet uitgewerkt, controleer met de afgeleide.)\label{vb:h3-vormenkort}}{%
\uitleg{Lees eerst welke vorm het is: de rest volgt dan vanzelf.}
\begin{tabular}{@{}*{2}{>{$\displaystyle}l<{$}@{\hspace{9mm}}}@{}}
\ig \cos(3x + 1)\dx = \tfrac13 \sin(3x + 1) + C & \ig 3x^2\,\e^{x^3}\dx = \e^{x^3} + C \\[2.5ex]
\ig 2x(x^2 + 1)^3\dx = \tfrac14 (x^2 + 1)^4 + C & \ig \dfrac{1}{25 + x^2}\dx = \tfrac15 \bgtan \dfrac{x}{5} + C \\[2.5ex]
\ig \dfrac{2x}{x^2 + 5}\dx = \ln(x^2 + 5) + C & \ig \dfrac{1}{\sqrt{25 - x^2}}\dx = \bgsin \dfrac{x}{5} + C
\end{tabular}}

\begin{oefeningen}
\oef{h3-vormen}
\end{oefeningen}

\subsection{Substitutie bij een bepaalde integraal}

Bij een bepaalde integraal moet je de grenzen meenemen, want die horen bij $x$ en niet bij $u$. Waarom mag dat? De substitutie vervormt het gebied, maar de oppervlakte blijft dezelfde. In het voorbeeld hieronder wordt $x$ van 0 tot 1 omgezet in $u$ van 1 tot 2. De linkse figuur toont $x\sqrt{1 + x^2}$ en de rechtse $\tfrac12\sqrt{u}$. De twee grafieken zien er verschillend uit, maar ze sluiten even veel oppervlakte in.

\figuren{\begin{tikzpicture}[x=1.6cm,y=1.35cm,font=\scriptsize]
  \gebied{0}{1}{\x*sqrt(1+(\x)^2)}
  \rooster{0}{1}{1}{0}{1}{1}
  \kromme{0}{1.05}{\x*sqrt(1+(\x)^2)}
\end{tikzpicture}\qquad\qquad
\begin{tikzpicture}[x=1.6cm,y=1.35cm,font=\scriptsize]
  \gebied{1}{2}{.5*sqrt(\x)}
  \rooster{0}{1}{2}{0}{1}{1}
  \kromme{.2}{2.05}{.5*sqrt(\x)}
\end{tikzpicture}}

\voorbeeld{}{Bereken $\ig_0^1 x\sqrt{1 + x^2}\dx$.\label{vb:h3-subgrens}}{%
\uitleg{Bij een bepaalde integraal zetten we ook de grenzen om. Zo moeten we op het einde niet meer terug naar $x$.}
\begin{stappen}
\text{Stel } u = 1 + x^2, \quad x\dx = \tfrac12\du & Grenzen: $x = 0$ geeft $u = 1$, en $x = 1$ geeft $u = 2$. \\[1ex]
\tfrac12 \ig_1^2 u^{1/2}\du = \tfrac12 \grens{\tfrac23 u^{3/2}}{1}{2} = \dfrac{2\sqrt{2} - 1}{3} & De integraal staat nu in $u$, met de nieuwe grenzen. $2^{3/2} = 2\sqrt{2}$ en $1^{3/2} = 1$.
\end{stappen}}

\voorbeeld{}{Bereken $\ig_1^{\e} \dfrac{(\ln x)^2}{x}\dx$.\label{vb:h3-subgrens2}}{%
\uitleg{Een breuk met $x$ in de noemer naast $\ln x$: de afgeleide van de binnenste functie staat er als factor $\tfrac1x$.}
\begin{stappen}
\text{Stel } u = \ln x, \quad \mathrm{d}u = \dfrac{1}{x}\dx & Grenzen: $x = 1$ geeft $u = 0$, en $x = \e$ geeft $u = 1$. \\[2ex]
\ig_0^1 u^2\du = \dfrac13 & De oude grenzen 1 en $\e$ zouden hier een fout geven.
\end{stappen}}

\begin{werkwijze}
Reken de grenzen meteen om naar $u$: $x = a$ wordt $u = g(a)$ en $x = b$ wordt $u = g(b)$. Daarna hoef je niet meer terug naar $x$.
\end{werkwijze}

\nujij{Bereken $\int_0^{\pi/2} \sin^2 x \cos x\dx$.}{$u = \sin x$, grenzen 0 en 1: $\int_0^1 u^2\du = \tfrac13$.}

\begin{oefeningen}
\oef{h3-subgrens}
\end{oefeningen}

\subsection{Eerst omvormen}

Soms staat er geen samengestelde functie met de juiste afgeleide. Bij $\sin^2 x$ ontbreekt de factor $\cos x$ om $u = \sin x$ te stellen. Een goniometrische formule helpt dan: de formules van Carnot, die je al kent, halen het kwadraat weg.
\[ \sin^2 x = \frac{1 - \cos 2x}{2} \qquad\text{en}\qquad \cos^2 x = \frac{1 + \cos 2x}{2} \]

\voorbeeld{}{Bereken $\ig_0^{\pi} \sin^2 x\dx$.\label{vb:h3-omvormen}}{%
\uitleg{De afgeleide van $\sin x$ ontbreekt, dus substitutie helpt niet. Met de formule van Carnot verdwijnt het kwadraat.}
\begin{stappen}
\ig_0^{\pi} \dfrac{1 - \cos 2x}{2}\dx = \grens{\dfrac{x}{2} - \dfrac{\sin 2x}{4}}{0}{\pi} = \dfrac{\pi}{2} & Het is niet $\dfrac{\sin^3 x}{3}$: de afgeleide van $\sin x$ ontbreekt.
\end{stappen}}

\naast{De uitkomst $\tfrac{\pi}{2}$ kun je ook aflezen. De grafiek van $\sin^2 x$ schommelt op $[0, \pi]$ tussen 0 en 1, en het stuk boven de rechte $y = \tfrac12$ past precies in het stuk eronder. De gemiddelde waarde is dus $\tfrac12$, en de integraal is $\tfrac12 \cdot \pi$.}{\begin{tikzpicture}[x=1cm,y=1.1cm,font=\scriptsize]
  \gebied{0}{3.1416}{(sin(\x r))^2}
  \roosterleeg{0}{1}{4}{0}{1}{1}
  \kromme{0}{3.1416}{(sin(\x r))^2}
  \draw[accent,dashed] (0,.5) -- (3.1416,.5);
  \node[below] at (3.1416,0) {$\pi$};
  \node[left] at (0,.5) {$\frac12$};
  \node[left] at (0,1) {1};
\end{tikzpicture}}

In andere gevallen blijft er na de substitutie nog een $x$ over. Die druk je dan ook uit in $u$.

\voorbeeld{}{Bereken $\ig x\sqrt{x + 1}\dx$.\label{vb:h3-rest}}{%
\uitleg{De wortel is het lastigste, dus die pakken we aan. Er blijft een $x$ over, maar die kunnen we uit de substitutie halen.}
\begin{stappen}
\text{Stel } u = x + 1, \quad x = u - 1, \quad \mathrm{d}u = \dx & In de wortel zit $x + 1$. De overgebleven $x$ schrijf je ook met $u$. \\[1ex]
\ig (u - 1)\,u^{1/2}\du = \ig (u^{3/2} - u^{1/2})\du & Uitwerken: $u \cdot u^{1/2} = u^{3/2}$. Nu zijn het enkel machten. \\[2.5ex]
= \tfrac25 (x+1)^{5/2} - \tfrac23 (x+1)^{3/2} + C & Terug naar $x$.
\end{stappen}}

\nujij{Bereken $\int \cos^2 x\dx$.}{$\tfrac{x}{2} + \tfrac{\sin 2x}{4} + C$.}

\begin{oefeningen}
\oef{h3-omvormen}
\end{oefeningen}

\subsection{Twee keer substitueren}

Soms is de integrand na één substitutie nog niet eenvoudig genoeg, maar herken je in de nieuwe integraal opnieuw een binnenste functie met haar afgeleide. Dan substitueer je gewoon nog een keer.

\voorbeeld{}{Bereken $\ig \dfrac{\sin^2(\ln x)\cos(\ln x)}{x}\dx$.\label{vb:h3-dubbel}}{%
\uitleg{De factor $\tfrac1x$ verraadt de afgeleide van $\ln x$, dus daarmee beginnen we. Daarna staat er een bekende vorm, waar nog een substitutie op past.}
\begin{stappen}
\text{Stel } u = \ln x, \quad \mathrm{d}u = \dfrac{1}{x}\dx & De eerste substitutie ruimt de logaritme op. \\[2ex]
\ig \sin^2 u\cos u\du & Nu is $\cos u$ de afgeleide van $\sin u$. \\[2ex]
\text{Stel } v = \sin u, \quad \mathrm{d}v = \cos u\du & De tweede substitutie, met een nieuwe letter. \\[1ex]
\ig v^2\,\mathrm{d}v = \dfrac{v^3}{3} + C = \dfrac{\sin^3(\ln x)}{3} + C & Vervang in omgekeerde volgorde: eerst $v$, dan $u$.
\end{stappen}
In één keer $v = \sin(\ln x)$ stellen kan ook, maar dat zie je meestal pas achteraf.}

\voorbeeld{}{Bereken $\ig_0^1 \dfrac{x\,\e^{x^2}}{1 + \e^{x^2}}\dx$.\label{vb:h3-dubbel2}}{%
\uitleg{Met grenzen werkt het net zo, en nu gaan de grenzen telkens mee. Daardoor hoeven we niet terug naar $x$.}
\begin{stappen}
\text{Stel } u = x^2, \quad x\dx = \tfrac12\du & Grenzen: 0 en 1 blijven 0 en 1. \\[1ex]
\tfrac12 \ig_0^1 \dfrac{\e^{u}}{1 + \e^{u}}\du & De afgeleide van de noemer staat nu in de teller. \\[2.5ex]
\text{Stel } v = 1 + \e^u, \quad \mathrm{d}v = \e^u\du & Grenzen: $u = 0$ geeft $v = 2$, en $u = 1$ geeft $v = 1 + \e$. \\[1ex]
\tfrac12 \ig_2^{1 + \e} \dfrac1v\,\mathrm{d}v = \tfrac12 \ln\dfrac{1 + \e}{2} & Numeriek ongeveer $0{,}310$.
\end{stappen}}

\begin{oefeningen}
\oef{h3-dubbel}
\end{oefeningen}

%========================================================================================
\section{Partiële integratie}

Substitutie helpt niet bij $x\,\e^x$: er is geen binnenste functie waarvan de afgeleide ernaast staat. Zo'n product van twee verschillende soorten functies vraagt een andere aanpak, en die komt van de productregel. Het idee is een ruil. Je leidt één factor af en integreert de andere, en krijgt een nieuwe integraal. Als die eenvoudiger is, heb je gewonnen.

\begin{theorie}
\begin{stelling}{Partiële integratie}
Voor een product met continue afgeleiden:

\vspace{1mm}
\qquad $\formule{\int f(x)\,g'(x)\dx = f(x)\,g(x) - \int f'(x)\,g(x)\dx}$

\vspace{1.5mm}
\qquad $\formule{\int_a^b f(x)\,g'(x)\dx = \Bigl[f(x)\,g(x)\Bigr]_a^b - \int_a^b f'(x)\,g(x)\dx}$
\end{stelling}

\begin{bewijs}
De productregel geeft
\bewijsregel{(fg)' = f'g + fg', \quad\text{dus}\quad fg' = (fg)' - f'g .}
Integreer beide leden. Omdat $fg$ een primitieve is van $(fg)'$, krijgen we
\bewijsregel{\ig fg'\dx = fg - \ig f'g\dx ,}
waarbij de constante in de overblijvende integraal zit. Met grenzen geeft de hoofdstelling
\bewijsregel{\ig_a^b fg'\dx = \Bigl[fg\Bigr]_a^b - \ig_a^b f'g\dx .}
\end{bewijs}
\end{theorie}

\naast{Je kunt de regel ook zien. Zet $g$ horizontaal en $f$ verticaal. De grote rechthoek min de kleine is $f(b)g(b) - f(a)g(a)$: dat is het stuk tussen rechte haken. De kromme verdeelt het gebied daartussen in twee delen. Het deel eronder is $\int f g'\dx$, het deel links ervan is $\int f' g\dx$. Samen vormen ze dus dat verschil. Ken je één deel, dan ken je ook het andere.}{\begin{tikzpicture}[x=1.15cm,y=1.05cm,font=\scriptsize]
  \fill[accent!30] (1,0) -- (1,1) .. controls (1.8,1.2) and (2.4,1.9) .. (3,2.6) -- (3,0) -- cycle;
  \fill[accent!13] (0,1) -- (1,1) .. controls (1.8,1.2) and (2.4,1.9) .. (3,2.6) -- (0,2.6) -- cycle;
  \draw[->,grijs] (-.1,0) -- (3.5,0) node[right,text=black] {$g$};
  \draw[->,grijs] (0,-.1) -- (0,3) node[above,text=black] {$f$};
  \draw[accent,thick] (1,1) .. controls (1.8,1.2) and (2.4,1.9) .. (3,2.6);
  \draw[lijngrijs,dashed] (1,0) -- (1,1) -- (0,1) (3,0) -- (3,2.6) -- (0,2.6);
  \node[below] at (1,0) {$g(a)$};
  \node[below] at (3,0) {$g(b)$};
  \node[left] at (0,1) {$f(a)$};
  \node[left] at (0,2.6) {$f(b)$};
  \node at (2.2,.6) {$\int f g'\dx$};
  \node at (1,1.9) {$\int f' g\dx$};
\end{tikzpicture}}

Welke factor moet je afleiden? Kies de factor die eenvoudiger wordt. Een veelterm daalt in graad, en een logaritme wordt een breuk. Een exponentiële functie blijft zichzelf, en een sinus of cosinus wordt niet eenvoudiger: die zijn dus een goede $g'$.

\voorbeeld{}{Bereken $\ig x\,\e^x\dx$.\label{vb:h3-partieel}}{%
\uitleg{Geen binnenste functie met de juiste afgeleide, dus partiële integratie. De keuze van $f$ en $g'$ is het belangrijkste: schrijf ze eerst op.}
\begin{stappenbreed}
f = x, \quad f' = 1, \quad g' = \e^x, \quad g = \e^x & We leiden $x$ af: dat geeft 1, en dan verdwijnt de $x$. Andersom zou er $x^2$ ontstaan. \\[1ex]
\ig x\,\e^x\dx = x\,\e^x - \ig \e^x\dx = x\,\e^x - \e^x + C & Vul in de formule in. Controle: $(x\e^x - \e^x)' = x\e^x$.
\end{stappenbreed}}

\voorbeeld{}{Bereken $\ig \ln x\dx$.\label{vb:h3-ln}}{%
\uitleg{Er lijkt maar één factor te staan, maar lees het als $\ln x \cdot 1$. De afgeleide van $\ln x$ kennen we wel, en haar primitieve (nog) niet.}
\begin{stappenbreed}
f = \ln x, \quad f' = \dfrac1x, \quad g' = 1, \quad g = x & We leiden de logaritme af. \\[2ex]
\ig \ln x\dx = x\ln x - \ig x \cdot \dfrac1x\dx = x\ln x - x + C & De $x$ valt weg, en er blijft een eenvoudige integraal over.
\end{stappenbreed}}

\begin{werkwijze}
Noem één factor $f$ en de andere $g'$. Kies zo dat $f'$ eenvoudiger wordt dan $f$, en dat je een primitieve $g$ kent. Een logaritme of een cyclometrische functie is meestal $f$. Staat er een veelterm naast $\e^x$, $\sin x$ of $\cos x$, dan is de veelterm $f$. Schrijf $f$, $f'$, $g'$ en $g$ eerst apart op.
\end{werkwijze}

\nujij{Bereken $\int x\cos x\dx$.}{$f = x$, $g' = \cos x$: $x\sin x - \int \sin x\dx = x\sin x + \cos x + C$.}

\begin{oefeningen}
\oef{h3-partieel}
\end{oefeningen}

\subsection{Twee keer, of rond in een cirkel}

Bij een hogere macht van $x$ herhaal je de methode: telkens daalt de graad met 1. Bij $\e^x$ met $\sin x$ of $\cos x$ komt na twee keer de oorspronkelijke integraal terug, en die los je op als een onbekende.

\voorbeeld{}{Bereken $\ig x^2 \cos x\dx$.\label{vb:h3-tweekeer}}{%
\uitleg{Hier volstaat één keer niet: de graad zakt telkens met 1, en pas bij graad 0 is de $x$ weg.}
\begin{stappenbreed}
\ig x^2 \cos x\dx = x^2 \sin x - \ig 2x \sin x\dx & $f = x^2$ en $g = \sin x$. De graad zakt van 2 naar 1. \\[2.5ex]
\ig 2x \sin x\dx = -2x\cos x + 2\sin x + C & Nog eens, met $f = 2x$ en $g = -\cos x$. Nu verdwijnt de $x$. \\[1ex]
\ig x^2 \cos x\dx = x^2 \sin x + 2x\cos x - 2\sin x + C & Let op het minteken voor de hele integraal.
\end{stappenbreed}}

\voorbeeld{}{Bereken $\ig \e^x \sin x\dx$.\label{vb:h3-rond}}{%
\uitleg{Geen van beide factoren wordt eenvoudiger. Na twee keer komen we terug bij de oorspronkelijke integraal, en dat is net de bedoeling.}
\begin{stappenbreed}
I = \ig \e^x \sin x\dx = \e^x \sin x - \ig \e^x \cos x\dx & $f = \sin x$ en $g = \e^x$. Geen van beide wordt eenvoudiger, maar dat geeft niet. \\[2.5ex]
\phantom{I} = \e^x \sin x - \e^x \cos x - I & Nog eens, met $f = \cos x$: de integraal keert terug. \\[1ex]
I = \tfrac12 \e^x (\sin x - \cos x) + C & Behandel $I$ als een onbekende: breng ze naar links, en deel door 2.
\end{stappenbreed}}

\nujij{Bereken $\int x^2 \e^x\dx$.}{$x^2\e^x - 2\int x\e^x\dx = \e^x(x^2 - 2x + 2) + C$.}

\begin{oefeningen}
\oef{h3-tweekeer}
\end{oefeningen}

\subsection{Partiële integratie met grenzen}

Met grenzen verandert er weinig. Het stuk tussen haken krijgt meteen zijn grenzen, en de integraal die overblijft ook. Je hoeft dus niet eerst de onbepaalde integraal op te lossen.

\voorbeeld{}{Bereken $\ig_1^{\e} x\ln x\dx$.\label{vb:h3-pgrens}}{%
\uitleg{De logaritme moeten we afleiden, want haar primitieve is niet eenvoudig. De rest gaat zoals zonder grenzen.}
\naast[]{\begin{stappenvrij}{7.1cm}{3.4cm}
f = \ln x, \quad g' = x, \quad g = \dfrac{x^2}{2} & De logaritme is $f$, dus $f' = \tfrac1x$. \\[2ex]
\grens{\dfrac{x^2}{2}\ln x}{1}{\e} - \ig_1^{\e} \dfrac{x}{2}\dx & Het stuk tussen haken krijgt al de grenzen. In de integraal valt $x$ weg tegen $\tfrac1x$. \\[3ex]
= \dfrac{\e^2}{2} - \grens{\dfrac{x^2}{4}}{1}{\e} = \dfrac{\e^2}{2} - \dfrac{\e^2 - 1}{4} = \dfrac{\e^2 + 1}{4} & $\ln 1 = 0$ en $\ln \e = 1$.
\end{stappenvrij}}{\begin{tikzpicture}[x=.95cm,y=.75cm,font=\scriptsize]
  \gebied{1}{2.718}{\x*ln(\x)}
  \rooster{0}{1}{3}{0}{1}{3}
  \kromme{.6}{2.9}{\x*ln(\x)}
  \node[above,fill=white,inner sep=.5pt] at (2.718,0) {$\e$};
\end{tikzpicture}}}

Numeriek is dat ongeveer 2,10. De figuur toont dat dit klopt: een gebied met basis 1,7 en een hoogte die oploopt tot 2,7.

\nujij{Bereken $\int_0^{\pi} x\cos x\dx$.}{$\bigl[x\sin x + \cos x\bigr]_0^{\pi} = -1 - 1 = -2$.}

\begin{oefeningen}
\oef{h3-pgrens}
\end{oefeningen}

%========================================================================================
\section{Welke methode?}

Dit is het moeilijkste van het hoofdstuk: de methode staat niet op de integraal geschreven. Integralen die erg op elkaar lijken, kunnen een andere aanpak vragen.

\voorbeeld{}{Vier integralen die op elkaar lijken: bereken $\ig x\,\e^{x^2}\dx$, \ $\ig x\,\e^{x}\dx$, \ $\ig (x + \e^{x})\dx$ en $\ig \dfrac{\e^{x}}{x}\dx$.\label{vb:h3-keuze}}{%
\uitleg{Alle vier bevatten ze $x$ en een exponentiële functie, maar de aanpak verschilt volledig. Kijk dus niet naar de ingrediënten, maar naar de bouw van de integrand.}
\begin{stappen}
\tfrac12 \e^{x^2} + C & Substitutie: $x$ is op een constante na de afgeleide van $x^2$. \\[1ex]
x\,\e^x - \e^x + C & Partiële integratie: geen binnenste functie, een product. \\[1ex]
\dfrac{x^2}{2} + \e^x + C & Splitsing: geen product maar een som. \\[1.5ex]
\text{geen primitieve in gekende functies} & De methodes lukken niet, en dat is geen tekort van jou.
\end{stappen}}

\naast{Dat laatste verdient aandacht. Veel functies hebben geen primitieve die je met gekende functies kunt schrijven, zoals $\e^{-x^2}$, $\tfrac{\sin x}{x}$ of $\sqrt{1 + x^3}$. Hun bepaalde integraal bestaat wel: de grafiek van $\e^{-x^2}$ sluit op $[0, 1]$ gewoon een gebied in. Een rekentoestel of computer benadert die oppervlakte met een som, zoals in hoofdstuk 1, en vindt $\int_0^1 \e^{-x^2}\dx \approx 0{,}747$.}{\begin{tikzpicture}[x=.85cm,y=1.1cm,font=\scriptsize]
  \gebied{0}{1}{exp(-(\x)^2)}
  \rooster{-2}{1}{2}{0}{1}{1}
  \kromme{-2}{2}{exp(-(\x)^2)}
\end{tikzpicture}}

\begin{werkwijze}
Doorloop deze vragen.
\begin{enumerate}[label=\arabic*.,itemsep=0pt,midpenalty=0]
  \item Staat de functie in de tabel, of na uitwerken of splitsen wel? Dan is het \emph{splitsing}.
  \item Staat er een samengestelde functie, en is de afgeleide van de binnenste functie ook aanwezig (op een constante na)? Dan is het \emph{substitutie}.
  \item Staat er een product van verschillende soorten functies, of een logaritme of cyclometrische functie alleen? Dan is het \emph{partiële integratie}.
  \item Een macht van $\sin x$ of $\cos x$? Een even macht: \emph{dubbele hoek}. Een oneven macht: één factor apart houden en \emph{substitutie}.
\end{enumerate}
Controleer het antwoord altijd met de afgeleide.
\end{werkwijze}

\begin{oefeningen}
\oef{h3-keuze}
\end{oefeningen}

%========================================================================================
\opdrachtenkop

\begin{oefeningen}
\needspace{24\baselineskip}
\oef{h3-d1}

\needspace{24\baselineskip}
\oef{h3-d2}

\needspace{9\baselineskip}
\oef{h3-d3}

\needspace{24\baselineskip}
\oef{h3-d4}

\needspace{17\baselineskip}
\oef{h3-d5}

\oef{h3-o14}

\oef{h3-oln}

\oef{h3-osc}

\oef{h3-o19}

\oef{h3-o20}

\oef{h3-o21}

\oef{h3-o23}

\needspace{20\baselineskip}
\oef{h3-d6}

\oef{h3-weierstrass}

\oef{h3-hyperbolisch}

\oef{h3-o24}

\oef{h3-o25}

\oef{h3-o26}

\oef{h3-o27}
\end{oefeningen}

%----------------------------------------------------------------------------------------
% KOPPEL: verbanden met andere leerstof
%----------------------------------------------------------------------------------------
\needspace{20\baselineskip}
\opdrachtsectie{Koppel}

\lkr{Complexe getallen en matrices zijn gezien in het 5de jaar. De oefening over Rijen gebruikt recursieve rijen, en die komen volgens de README van het pakket pas na Analyse 3: kijk dat na. De vergelijking van een ellips staat bij de suggesties voor het complementair gedeelte, maar wordt in de oefening gegeven.}

\begin{oefeningen}
\oef{h3-koppel-1}

\oef{h3-koppel-2}

\oef{h3-koppel-3}

\oef{h3-koppel-4}
\end{oefeningen}

%----------------------------------------------------------------------------------------
\opdrachtsectie{Test jezelf}

\begin{oefeningen}
\oef{h3-t1}

\oef{h3-t2}

\oef{h3-t3}

\oef{h3-t5}

\oef{h3-t7}

\oef{h3-t8}

\oef{h3-t10}
\end{oefeningen}

\herhaalschema{%
  \fv{oef:h3-t1}{vb:h3-voorwaarde}{oef:h3-onb}%
  \fv{oef:h3-t2}{vb:h3-split}{oef:h3-split}%
  \fv{oef:h3-t3}{vb:h3-vormen}{oef:h3-vormen}%
  \fv{oef:h3-t5}{vb:h3-subgrens}{oef:h3-subgrens}%
  \fv{oef:h3-t7}{vb:h3-pgrens}{oef:h3-pgrens}%
  \fv{oef:h3-t8}{vb:h3-rond}{oef:h3-tweekeer}%
  \fv{oef:h3-t10}{vb:h3-keuze}{oef:h3-keuze}}

\ifzelfstudie
\section*{Maak je eigen samenvatting}
Sluit dit boek en schrijf in je schrift: wat een onbepaalde integraal is, de regel voor substitutie met zijn bewijs, de vormen die vaak terugkomen, de regel voor partiële integratie met zijn bewijs, hoe je $f$ en $g'$ kiest, en de vier vragen om een methode te kiezen. Vergelijk daarna met de kaders.
\fi
